%%%%%%%%%%%%%%%%%%%%%%%%%%%%%%%%%%%%%%%%%%%%%%%%%%%%%%%%%%%%%%%%%
%% Übersichtstabelle für Dateien und Verzeichnisse der Vorlage %%
%%%%%%%%%%%%%%%%%%%%%%%%%%%%%%%%%%%%%%%%%%%%%%%%%%%%%%%%%%%%%%%%%
%% Spaltendefinition:
%% Erste Spalte vom Typ "l": Spalte so breit, wie der Inhalt, linksbündige Ausrichtung
%% Zweite Spalte vom Typ "x": die Breite wird automatisch berechnet
%% in Abhängigkeit von der Gesamtbreite der Tabelle und Breite anderer Spalten,
%% Inhalt wird umbrochen, wenn zu lang
%% Format: linksbündiger Flattersatz ohne Silbentrennung, 
%% Weitere Möglichkeiten:
%% "X" = Blocksatz (mit Silbentrennung)
%% "F" = linksbündiger Flattersatz mit Silbentrennung (KSP-Verlag: unerwünscht),
%% "y" = rechtsbündiger Flattersatz ohne Silbertrennung
%% "Y" = rechtsbündiger Flattersatz mit Silbertrennung
%% Spacing zwischen den Spalten auf 8pt setzen:
\begin{longtable}{l@{\extracolsep{8pt}}x}%
%-----------------------------------------------------------------------------------------------------
\caption[Dateien und Verzeichnisse der Vorlage]{Dateien und Verzeichnisse der Vorlage}%
\label{tab:StrukturDerVorlage}%
\tabularnewline%
%-----------------------------------------------------------------------------------------------------
\toprule%
%-----------------------------------------------------------------------------------------------------
\bfseries Datei/Verzeichnis               & \bfseries Bedeutung und Benutzerinteraktion
\tabularnewline%
%-----------------------------------------------------------------------------------------------------
\midrule%
\endfirsthead%
%-----------------------------------------------------------------------------------------------------
\toprule%
%-----------------------------------------------------------------------------------------------------
\bfseries Datei/Verzeichnis               & \bfseries Bedeutung und Benutzerinteraktion
\tabularnewline%
%-----------------------------------------------------------------------------------------------------
\midrule%
\endhead%
%-----------------------------------------------------------------------------------------------------
\printfilepath{./Diss.tex}                       & \texttt{TeX}-Hauptdatei mit Einbindung weiterer Dateien.
																								 Hier kann die Kapitalstruktur angepasst werden.\\
%-----------------------------------------------------------------------------------------------------
\printfilepath{./.gitignore}                     & Spezifikation von Dateien und Verzeichnissen, die von der Versionskontrolle ignoriert werden sollen.
																								 Normalerweise keine Anpassung nötig.\\
%-----------------------------------------------------------------------------------------------------
\printfilepath{./.latexmkrc}                     & Konfigurationsdatei für den Aufruf von \texttt{latexmk}\index{latexmk}.
																								 Normalerweise keine Anpassung nötig.\\
%-----------------------------------------------------------------------------------------------------
\printfilepath{./bib/Diss.bib}                   & Literaturdatenbank im \texttt{BibLaTeX}\index{biblatex}-Format.
																								 Hier verwendete Referenzen einfügen.\\
%-----------------------------------------------------------------------------------------------------
\printfilepath{./content/*}                      & Eingebundene Inhalte der Arbeit.
																								 Hier Inhalte der einzelnen Kapitel verwalten.\\
%-----------------------------------------------------------------------------------------------------
\printfilepath{./figures-src/*}                  & \texttt{TikZ}\index{TikZ}\index{Bilder!TikZ}-Zeichnungen. %(Quellcode).
                                                 Ggf. weitere hinzufügen.\\
%-----------------------------------------------------------------------------------------------------
\printfilepath{./figures-compiled/*}             & Temporäre Kompilate der \texttt{TikZ}-Zeichnungen.
                                                 Bei Aktualisierung der \texttt{TikZ}-Zeichnungen löschen.\\
%-----------------------------------------------------------------------------------------------------
\printfilepath{./fonts/*}                        & Verwendete Schriften.
																								 Normalerweise keine Anpassung nötig.\\
%-----------------------------------------------------------------------------------------------------
\printfilepath{./images/*}                       & Bilder\index{Bilder} im Binärformat (jpg, png, pdf, \etc).
                                                 Ggf. eigene Bilder hinzufügen. \\
%-----------------------------------------------------------------------------------------------------
\printfilepath{./preambel/*}                     & Konfigurationsdateien (s. u.)\\
%-----------------------------------------------------------------------------------------------------
\printfilepath{./preambel/Acronyms.tex}          & Definition der Akronyme.
                                                 Ggf. eigne hinzufügen.\\
%-----------------------------------------------------------------------------------------------------
\printfilepath{./preambel/AlleAngaben.tex}       & Wichtige Angaben und Einstellungen.
                                                 Hier Angaben zum Typ der Arbeit, zum Autor und zu den Gutachtern ändern.\\
%-----------------------------------------------------------------------------------------------------
\printfilepath{./preambel/AllePfade.tex}         & Definition von Suchpfaden für Bildverzeichnisse und Bibliografie"=Dateien.
                                                 Ggf. ergänzen.\\
%-----------------------------------------------------------------------------------------------------
\printfilepath{./preambel/AlleSchalter.tex}      & Schalter zur Umstellung der Hauptsprache des Manuskriptes,
                                                 zur Einstellung der Druckfarben sowie zur
                                                 Aus- und Wieder"=Einblenden der einzelnen Manuskript"=Teile.
                                                 Bei Bedarf umstellen.\\
%-----------------------------------------------------------------------------------------------------
%\printfilepath{./preambel/BibSettings.tex}       & Konfiguration der Bibliografie
%                                                 Normalerweise keine, es sei denn man möchte den Stil der Literaturverzeichnisse ändern.\\
%-----------------------------------------------------------------------------------------------------
\printfilepath{./preambel/EncodingAndFont.tex}   & Schriftarteinstellungen.
                                                 Keine Anpassung nötig, sofern man die vorgegeben Schriftarten nutzen möchte.\\
%-----------------------------------------------------------------------------------------------------
\printfilepath{./preambel/Glossary.tex}          & Definition der Glossar-Einträge. 
                                                 Hier können ggf. weitere Einträge hinzugefügt werden.\\
%-----------------------------------------------------------------------------------------------------
\printfilepath{./preambel/GlossarySymbols.tex}   & Definition der Glossar-Einträge für die automatische Erzeugung des Symbolverzeichnisses.
                                                 Bei Verwendung um eigene Symboldefinitionen ergänzen.\\
%-----------------------------------------------------------------------------------------------------
\printfilepath{./preambel/Header.tex}            & Alle Präambel-Definitionen (\teilw in weiteren Dateien).
                                                 Aktivierung der \printkeyword{draft}-Option und des A4-Layouts.\\
%-----------------------------------------------------------------------------------------------------
\printfilepath{./preambel/Hyphenation.tex}       & Silbentrennung für unbekannte Wörter.
                                                 Hier können ggf. Regeln für die Silbentrennung weiterer Begriffe hinzugefügt werden.\\
%-----------------------------------------------------------------------------------------------------
\printfilepath{./preambel/IndexStyle.tex}        & Layout des Stichwortverzeichnisses.
                                                 Keine Anpassung nötig.\\
%-----------------------------------------------------------------------------------------------------
\printfilepath{./preambel/KomaOptions.tex}       & KomaScript-Optionen.
                                                 Keine Anpassung nötig.\\
%-----------------------------------------------------------------------------------------------------
\printfilepath{./preambel/Math.tex}              & Mathe-Einstellungen und Makros.
                                                 Bei Bedarf eigene Mathe-Makros definieren.\\
%-----------------------------------------------------------------------------------------------------
\printfilepath{./preambel/MyPackages.tex}        & Eingebundene \LaTeX"=Zusatzpakete.
                                                 Bei Bedarf können hier Zusatzpakete eingebunden werden.\\
%-----------------------------------------------------------------------------------------------------
\printfilepath{./preambel/Newcommands.tex}       & Eigene \LaTeX{}-Makros.
                                                 Hier können ggf. eigene Makros hinzugefügt werden.\\
%-----------------------------------------------------------------------------------------------------
\printfilepath{./preambel/preambel-commands.tex} & Interne Befehle aus der alten Vorlage.
                                                 Normalerweise keine Anpassung nötig.\\
%-----------------------------------------------------------------------------------------------------
\printfilepath{./preambel/settings.tex}          & Einstellungen zu Längen, Breiten, Verzeichnistiefen etc.
                                                 Normalerweise müssen keine Änderungen vorgenommen werden, 
                                                 da diese Einstellungen mit dem \gls{ac:KSP} abgestimmt 
                                                 und bereits in der Vorlage umgesetzt wurden.\\
%-----------------------------------------------------------------------------------------------------
\printfilepath{./preambel/TableCommands.tex}     & Tabelleneinstellungen.
                                                 Normalerweise keine Anpassung nötig.\\
%-----------------------------------------------------------------------------------------------------
\printfilepath{./preamble/Translations.tex}      & Multilinguale Begriffsdefinitionen für Beschriftungen etc.
																				         Normalerweise keine Anpassung nötig.\\
%-----------------------------------------------------------------------------------------------------
\printfilepath{./tables/*}                       & In eigene Dateien ausgelagerte Tabellen.
                                                 Ggf. eigene hinzufügen. \\
%-----------------------------------------------------------------------------------------------------
\bottomrule%
%-----------------------------------------------------------------------------------------------------
\end{longtable}%